\documentclass[10pt]{article}
\usepackage[margin=20mm]{geometry}
\usepackage{amsmath,amssymb,mathtools,fontspec,unicode-math,xurl,hyperref}
\IfFontExistsTF{DejaVu Serif}{\setmainfont{DejaVu Serif}}{\setmainfont{Latin Modern Roman}}
\setmathfont{latinmodern-math.otf}
\setlength{\parindent}{0pt}\setlength{\parskip}{6pt}\sloppy
\title{Central terms, degree-zero conventions, and explicit formulas in the spin GW/Hurwitz correspondence}
\author{AI-assisted technical review}
\date{}
\begin{document}
\maketitle
Abstract: The 35-page v3 manuscript of Giacchetto, Kramer, Lewański, and Sauvaget and its nine-page computational supplement are examined. Finite calculations expose incompatible empty-partition and degree-zero conventions, an omitted central term in the Hodge operator, and an incorrect one-point recurrence. The printed degree-zero geometric formula also disagrees with the stated localization series. The connectedness of the principal statement and the positive-form proof of the conditional extension also require substantive correction. Local corrections concern operator signs and the scope of a Lie algebra projection. The supplement’s low-degree formulas and corrected recurrence receive independent exact checks. These findings do not establish that the intended positive-degree correspondence or BKP conclusion is false after repair. Scientific adjudication remains pending.\par
\section*{Identity, reading scope, and dependency structure}
The reviewed paper is A. Giacchetto, R. Kramer, D. Lewański, and A. Sauvaget, The spin Gromov–Witten/Hurwitz correspondence for \(\mathbb P^1\), DOI 10.4171/JEMS/1588. The primary scientific source is arXiv:2208.03259v3, 14 August 2025 [paper]. The initial review read the complete 35-page text and actually viewed every source page; it also read and viewed all nine pages of the author-linked Spin GW invariants of \(\mathbb P^1\): Calculations [supp]. Those retained facts are bound to unchanged source bytes. The present independent acceptance checks the eight reported Problems and their actual proof dependencies. Page references below refer to the primary PDF unless identified as supplemental.\par
The journal records acceptance on 8 October 2024 and publication on 6 January 2025 [journal]. Version 2, dated 17 October 2024, is labelled accepted on arXiv. Version 3 is later than publication; its authors report a correction to Lemma 2.11 and its proof. The v3 factor \(2^{-g}\) is actually present and is used in all calculations here. Its nonempty small cases are supported, not criticized. The retained v2 received only a targeted auxiliary comparison, not a second complete scientific review. The subscription journal PDF was not independently compared. The supplement has no independently established release date; its internal PDF metadata is not a publication date. Its older notation is not assumed identical to v3. No article-specific author code was found in the bounded acquired-source search. No author code was executed.\par
The proof route is: the neutral Clifford algebra and spin characters in Section 2; the spin ELSV theorem quoted there; the specialized localization calculation of Section 3; Hodge operators and their commutator in Section 4; the Fock-space GW formulas of Section 5; and dressing into the 2-BKP hierarchy in Section 6. Theorem E is explicitly conditional on Conjecture D, the proposed degeneration statement. It is not an unconditional theorem for every spin target curve. The cited spin ELSV, character theory, and hierarchy theorems have not been independently reproved here.\par
For clarity, suppress the superscript \(B\) on operators and set \[A=\alpha_1,\quad B=\alpha_{-1},\quad
 s(z)=2\sinh(z/2),\quad q_0(z)=\tfrac12\cosh(z/2),\quad
 c(z)=\frac{q_0(z)}{s(z)}=\tfrac14\coth(z/2).\] Thus \([A,B]=1/2\) and \(\mathcal E_0=\widehat{\mathcal E}_0+c\,\operatorname{Id}\). These definitions distinguish a scalar correction from a change in the nonzero-energy operators.\par
\section*{Substantive normalization and boundary issues}
Problem 1 (Major: completed cycles and removal of degree-zero terms are incompatible).\par
Location. Notation (p. 8), Definition 2.10 and (2.23) (p. 11), (2.24)–(2.27) (p. 12), Proposition 2.14, (2.31) (p. 13), the proof of Theorem C (p. 29), and Remark 5.7, (5.20) (p. 31).\par
The stated set of odd partitions includes \(\varnothing\). Equation (2.24) sums over it, and Lemma 2.11 gives \[\kappa_{k,\mu}
 =2^{-g}(k-1)!\frac{\prod_i\mu_i}{|\mu|!}
 [z^{2g}]S(2z)S(z)^{|\mu|-2}\prod_iS(\mu_i z),
 \qquad 2g=k+1-|\mu|-\ell(\mu).\] Since \(S(2z)/S(z)^2=(z/2)\coth(z/2)\), one obtains \[\kappa_{1,\varnothing}=\frac1{24},\quad
 \kappa_{1,(1)}=1,\qquad
 \kappa_{3,\varnothing}=-\frac1{1440},\quad
 \kappa_{3,(1)}=\frac16.\] In degree one, (2.23) gives \(H_1(\mathbb P^1;\varnothing)=H_1(\mathbb P^1;(1))=1\): the added-one binomial factor is one, and the unique degree-one cover has positive spin sign and trivial automorphism group. Consequently the literal completed-cycle expansion gives \[H_1(\overline c_1)=\frac{25}{24},\qquad
 H_1(\overline c_3)=\frac{239}{1440}.\] In contrast, the elementary Fock calculation \[\left\langle A\widehat{\mathcal E}_0(z)B\right\rangle=\tfrac12\sinh z\] makes (2.31) give \(1\) and \(1/6\), respectively. This contradiction already occurs in positive covering degree; it does not require a disputed genus convention for an empty source.\par
The proof (2.27) also uses an ordinary vacuum expectation where connected or inactive-sheet subtractions need to be specified. For example, with \(\mu=(1),k=1\) its ordinary expectation includes the same \(1/24\) vacuum contribution, although the final coefficient is \(1\). Changing \(\mathcal E_0\) to \(\widehat{\mathcal E}_0\) alone is not a proof of the required connected-cover extraction in general.\par
The proposed later repair for degree-zero components is itself incorrect. At \(n=1\), the left side of (5.20) is exactly \[c(z)\left\langle A^dB^d\right\rangle=\frac{d!}{2^d}c(z),\] whereas its right side is \(d!\,U_0(z)\). The cases \(d=1,2\) require the same \(U_0\) to equal both \(c/2\) and \(c/4\). For arbitrary \(n\), the actual difference of the two products is \[\sum_{S\subsetneq[n]}
 \left\langle A^d\prod_{i\in S}\widehat{\mathcal E}_0(z_i)B^d\right\rangle
 \prod_{i\notin S}c(z_i).\] It contains mixed positive-degree and scalar contributions.\par
Impact and repair boundary. The comparison used to prove Theorem C needs one consistent definition of completed cycles, connected extraction, and allowed degree-zero components. Excluding the empty partition agrees with the centered operator in these tests; retaining it requires the corresponding scalar and disconnected terms. Neither choice alone is claimed here to repair the entire correspondence. Proposition 5.4 already explicitly excludes degree-zero components and is not refuted by this calculation. The new v3 factor is not the source of this discrepancy.\par
Problem 2 (Major: the stable degree-zero geometric boundary is not covered consistently).\par
Location. (1.5)–(1.6) (p. 3), Definition 3.1 and its unstable conventions (pp. 14–15), (3.14)–(3.16), Theorem 3.2 (pp. 16–18).\par
Theorem 3.2 includes \(d=0\), and (1.5) imposes no positive-degree restriction on its Chern-class formula. For the stable case \(g=1,d=0,n=1\), \[\overline M_{1,1}(\mathbb P^1,0)=\overline M_{1,1}\times\mathbb P^1,
 \qquad c_{g-1+d}=c_0=1.\] Writing \(\lambda=c_1(\mathbb E)\) and \(h=[\mathrm{pt}]\), the ordinary virtual Euler class is \(e(\mathbb E^\vee\otimes T\mathbb P^1)=2h-\lambda\). The usual \(\int_{\overline M_{1,1}}\lambda=1/24\) therefore gives, from (1.5) as printed, \[\left\langle \tau_0(h)\right\rangle_{1,0}=-\frac1{24}.\] Its Definition 3.1 normalized coefficient is \(-z/12\). On the other hand, the empty-partition term of Theorem 3.2 uses \[H_1^\circ(z)
 =2\int_{\overline M_{1,1}}
 (1+2\lambda)(1-\lambda)\frac{z}{1-z\psi}
 =\frac{z+z^2}{12}.\] It gives \(+z/12\) in the nonequivariant limit. The discrepancy is in a stable case, so the separately assigned unstable one- and two-point terms do not settle it.\par
Connected positive-degree maps satisfy the required concavity; constant maps have \(H^0\neq0\) and \(R^1\) of rank \(g\), not \(g+d-1\). Restrict the geometric formula to its justified range and derive the degree-zero theory separately, or justify an explicit formal convention. The series alone does not establish the actual cosection-localized degree-zero value.\par
Independently, in the general \(\mathcal O(-k)\) formula (3.14), take a single degree-one edge with \(k=2,\alpha=1\). Its endpoint weights \([1,3]\) give \(\omega_{\mathbb P^1}\otimes\chi^2\), hence one-dimensional \(H^1\) of weight \(2\) and Euler class \(2t\). The printed product instead gives \[(1+2)(-t)+(1-2)t=-4t.\] The endpoint vertex factors are one. This does not refute the subsequent specialization \(k=1,\alpha=-1/2\): odd-edge signs cancel and even edges vanish. That specialized positive-degree calculation remains a separate question.\par
Problem 3 (Major: the Hodge operator omits a central term).\par
Location. Definition 4.1, (4.2), Proposition 4.2 (p. 20), its proof (pp. 21–22), and Theorem 4.7 (pp. 22–23).\par
The bound \(j\geq k'+1=(k+1)/2\) excludes \(k=j=0\). But \(\mathcal E_0(w)\) has leading scalar term \(1/(2w)\). At this index the Gamma quotient is \(1/z\), so the omitted contribution is the entire formal scalar \[\boxed{\frac{e^{uz^3/12}}{2z}\operatorname{Id}.}\] In particular the printed definition gives \(\mathcal B(1,0)=B\), whereas (4.3) gives \[e^A B e^{-A}=B+\tfrac12\operatorname{Id}.\] Theorem 4.7 at \(n=1,u^0\) likewise requires the nonzero unstable Hodge term \(1/(2z)\).\par
The printed \(u^1\) vacuum coefficient \(z(z+2)/24\) gains \(z^2/24\) from the missing scalar, giving \(H_1^\circ(z)=z(z+1)/12\). Restoring only \(1/(2z)\) therefore fails already at the next coefficient.\par
Impact and repair boundary. Scalar terms alter Hodge/GW expectations, not commutators. The proof on p. 22 drops \(k=0\). Restore the full scalar throughout; the two checked coefficients do not certify all-genus or central-lift identities.\par
\section*{Statement errors and further proof restrictions}
Problem 4 (Major: dimension and connectedness in the principal statements).\par
Location. (1.7), Theorem A (p. 3), Conjecture B (p. 4), Notation (p. 8), and the proofs on pp. 28–29.\par
Conjecture B prints \(\sum k_i=g-1+d(1-g(B))+n\). Each stationary insertion has codimension \(k_i+1\), so the \(n\) cancels against the localized virtual dimension. For \(\mathbb P^1\) the condition is \(\sum k_i=g-1+d\), precisely the condition actually used on p. 29. For example \(g=0,d=1,n=1,k_1=1\) satisfies the printed condition, but \(\psi_1\operatorname{ev}_1^*h\) has codimension two on the one-dimensional degree-one moduli space and integrates to zero. Using the literal completed-cycle convention of Problem 1, the stated corresponding expression is \(-\tfrac12H_1(\overline c_3)=-239/2880\neq0\). This example keeps that convention fixed throughout. Even using the centered coefficient \(1/6\) instead would give \(-1/12\), so the dimension conflict does not depend on retaining Problem 1’s empty-cycle contribution.\par
Equation (1.7) has no disconnected marker, although the notation declares an unmarked bracket connected. Set \(x=x^\star=0\). The connected \(d=2,g\geq0\) integrals without insertions vanish by their positive dimension \(g+1\). Theorem A’s operator expression instead equals \[\left\langle e^A(q/u)^H e^B\right\rangle=\exp(q/(2u)),\] whose \(q^2\) coefficient is \(1/(8u^2)\). An empty-curve convention can change the constant term but not this coefficient. The power \(u^{-2}\) cannot occur in a connected genus sum with \(g\geq0\); it represents two disconnected degree-one rational components. This refutes the literal principal statement without using any problematic inserted operator from Problems 1–3. Theorem 5.3 already has the proper bullet, which identifies the intended repair but does not validate Theorem A as stated. Use the disconnected potential with its arithmetic-genus range and delete the extra \(n\) from the dimension condition. The severity is Major because the asserted main result, not only auxiliary notation, changes.\par
Problem 5 (Minor: two operator signs disagree with their actual use).\par
Location. (2.12)–(2.13) (p. 10), (2.19)–(2.20) (p. 11), (2.30) (p. 13), and the proofs on pp. 21, 23.\par
The exponent in (2.12)–(2.13) is \(e^{(k+m/2)z}\). With the displayed indices \(k-m,k\), the centered exponent needed by the subsequent parity relation is \(e^{(k-m/2)z}\). Directly from the CAR, \[\mathcal E_{-1}^{\rm printed}(z)|0\rangle
 =e^{-z}\cosh(z/2)\phi_1\phi_0|0\rangle,\qquad
 \left\langle A\mathcal E_{-1}^{\rm printed}(z)\right\rangle
 =\tfrac12e^{-z}\cosh(z/2).\] Its derivative at zero is \(-1/2\), whereas (2.19) requires it to be even. Correcting the sign fixes this finite test; the proof on p. 21 already uses \(e^{w(l-k/2)}\).\par
The right exponential in each conjugated factor of (2.30) also needs a minus sign. For \(r=2,\mu=(1,1)\), put \(v_j=\phi_j\phi_0|0\rangle\). Then \(B|0\rangle=v_1\), \(Bv_1=v_2\), \(F_3v_j=j^3v_j\), and \(\left\langle e^A v_2\right\rangle=1/4\). Formula (2.29) has vacuum value \(e^{8u/3}/4\); printed (2.30) has \(e^{10u/3}/4\). With \(g=0,b=1\) the numbers are \(4/3\) and \(5/3\). The proof on p. 23 already uses the correct inverse exponential. Both findings are definition/use mismatches, with evidence of the intended local repair.\par
Problem 6 (Minor: the full parity projection is not a Lie homomorphism).\par
Location. Proposition 4.8 and its proof (pp. 23–24); the subsequent dressing argument (pp. 24–26).\par
Even after correcting the centering, the A-type relation gives \[[\mathcal E_2(0),\mathcal E_0(w)]=s(2w)\mathcal E_2(w),\qquad
 [\mathcal E_{2,0},\mathcal E_{0,2}]=2\mathcal E_{2,1}.\] The proposed \(p\) kills both inputs, since their index-plus-degree is even. It sends the bracket to \(2\mathcal E^B_{2,1}\neq0\). For a finite nonzero test, the vacuum component of \(\mathcal E^B_2(w)v_2\) is \(\tfrac12\sinh(w)|0\rangle\). Thus the coefficient of \(w\) does not vanish.\par
A narrower algebraic route is available. Let \(\tau\) be the anti-automorphism with \(\tau(S)=-S\), \(\tau(H)=-H\). The \(\tau\)-anti-invariant elements form a Lie subalgebra; the odd parity sector embeds correctly there. The actual operators \(D,\widetilde D,X,E\) of Lemma 4.9 are anti-invariant. For \(W'=WE/t\), \(W(0)=1\), one obtains \(\tau(W)=W^{-1}\) whenever this formal solution is defined, so its conjugation preserves the required subalgebra. This supplies a restricted route for the displayed calculations, not the asserted homomorphism from the full algebra. The choices of formal completion, central lift, and triangular normalization still require explanation. In particular the actual definition is \(Z=tS^2/(2u)\); it has not been silently replaced by its reciprocal. No counterexample to the commutator formula itself is asserted.\par
Problem 7 (Minor: the one-point recurrence and auxiliary energy are wrong).\par
Location. Corollary 5.5, (5.14)–(5.16) (p. 30); supplement, pp. 1–2, Corollary 2, Definition 4, Lemma 5.\par
With \(s=s(z),q_0=q_0(z)\), the printed recurrence starts at \((U_1,V_1)=(sq_0,q_0)\) and gives \[U_2=sq_0(s^2+4)=\tfrac12\sinh(2z)+\sinh z.\] The degree-two Fock space has the single strict partition \((2)\) with norm weight \(1/2\), hence \(U_2=\tfrac12\sinh(2z)\). For the valid parameters \(d=2,n=1,k=1,g=0\), (5.13) therefore gives \(-5/12\) from the printed recurrence, but \(-1/3\) from Corollary 5.6 and the actual operator expression.\par
Equation (5.16) defines \(V_d\) with total nonzero energy: \(\left\langle A^d\mathcal E_1(z)B^{d-1}\right\rangle=0\), contradicting \(V_1=q_0\). The supplement supplies the correctly balanced definition \(V_d=\left\langle A^{d-1}\mathcal E_1(z)B^d\right\rangle\) and the relations \[\begin{split}
 V_d&=sU_{d-1}+\frac{d-1}{2}V_{d-1}
       +\frac{(d-1)!}{2^{d-1}}q_0,\\
 U_d&=sV_d+\frac d2 U_{d-1}.
 \end{split}\] They follow by commuting \(A\) and \(B\), using \([A,B]=1/2\), \([A,\widehat{\mathcal E}_0]=s\mathcal E_1\), and \([\mathcal E_1,B]=s\widehat{\mathcal E}_0+q_0\). The corrected matrix and inhomogeneous term are thus \[\begin{pmatrix}s^2+d/2&(d-1)s/2\\s&(d-1)/2\end{pmatrix},
 \qquad
 \frac{(d-1)!}{2^{d-1}}q_0\binom{s}{1}.\] The initial review’s retained finite occupation-space implementation of the CAR verified these recurrences and all eight rows of the supplement’s Table 1 for \(1\leq d\leq8\). It also recovers the multipoint spectral sums \[\begin{split}
 d=1 &: \tfrac12\prod_i\sinh z_i,\\
 d=2 &: \tfrac12\prod_i\sinh(2z_i),\\
 d=3 &: \tfrac12\prod_i\sinh(3z_i)
       +\tfrac14\prod_i(\sinh(2z_i)+\sinh z_i).
 \end{split}\] These are counterevidence to any claim that all of the low-degree formulas are wrong. The recurrence can be corrected without discarding Proposition 5.4.\par
Problem 8 (Major: the conditional extension uses an undefined positive form).\par
Location. Section 5.2, (5.11)–(5.12) and the following conclusion (p. 29).\par
Under the usual evaluation of supersymmetric functions on strict partitions, bounded polynomial degree does not give finite support. For \(\phi=\psi=1,\xi=1\), the terms with \(\lambda=(r)\) in \[q_\xi(\phi,\psi)=
 \sum_{\lambda\in SP}\xi^{|\lambda|}2^{-p(\lambda)}
 \dim(V^\lambda)\phi(\lambda)\psi(\lambda)\] are each at least \(1/2\). The sum diverges. Thus the asserted finiteness and convergence for every \(\xi\) fail under that interpretation. Reading \(\phi(\lambda)\) as Schur-\(Q\) coordinates would give a different, finitely supported meaning. That meaning is not defined here, and (5.12) must be proved for it. In particular, (2.23) extends the empty partition as an identity in every degree; the empty coordinate is supported only in degree zero. These interpretations cannot be interchanged without proving compatibility with the Hurwitz pairing.\par
A possible proof would use a sufficiently large finite partition cutoff on each filtered finite-dimensional subspace, prove nondegeneracy of the resulting positive form, and justify (5.12) coefficient by coefficient. A unipotent isometry of a positive finite-dimensional form is the identity. The missing separation and identification steps have not been supplied here. The valid finite-dimensional isometry argument does not supply the missing Gram form or identity (5.12). These are the actual steps used to identify the two completed-cycle bases and prove Theorem E from Conjecture D. Their failure is therefore a Major gap in that conditional proof; no counterexample to Theorem E or Conjecture D is claimed.\par
\section*{Coverage, checks, and assessment}
The localization parity filter and its specialized positive-degree edge factors, the stable Hodge vertex expressions, and the global sign and parameter bookkeeping were traced with their hypotheses intact. The decomposition of the identity in Lemma 5.2 and the stated exclusion of degree-zero components in Proposition 5.4 were retained. The Section 6 divisor and string arguments and the BKP dependency on the dressing operators were read, without treating unproved central/degree-zero consistency as automatic. No independent counterexample to a fully corrected positive-degree correspondence or BKP theorem has been produced.\par
The initial-review executor’s retained Python/SymPy run exited zero. Its actual assertions check the completed-cycle coefficients, Hodge scalar, parity test, CAR recurrence through degree eight, and strict-partition norm weights. The present executor inspected the unchanged script and output and independently checked the analytic arguments. The conjugation numbers in Problem 5 and bracket in Problem 6 are only printed constants in that script; their justification is the explicit operator calculation above. The genus-one geometry in Problem 2 is likewise analytic. No author code or moduli-space computational package was run.\par
The original native cross-checks and source viewing belong to the initial review. A native read-only auxiliary also challenged these eight Problems during acceptance and checked the exact supplement and calculation bindings. Native assistance is neither human review nor independent backend verification.\par
Assessment: major revision required. Problems 1–4 and 8 are Major: they concern core normalization, a stable boundary formula, the Hodge operator, literal principal statements, and the conditional extension’s proof. Problems 5–7 remain Minor, with explicit local repairs and counterevidence from the actual downstream formulas. A possible repair does not reduce the severity of a core invalid statement or inference. No claim is made that every principal conclusion is false after consistent repair. Scientific adjudication is pending; no human reviews are recorded.\par
\section*{Supporting checks}
The primary PDF has 599,803 bytes and SHA256\par
> f080b99ef6f1566aac35ee0c00508b7a3a205946d941c6a4d8c7d35852e429bc.\par
The supplement has 225,712 bytes and SHA256\par
> 6ace5bf2db15e90202e26327f860ee8234c1f70f893ded38bea88f7cfb96a50d.\par
The retained exact-check script has SHA256\par
> 0a39cf81eaec7a7950e21215e30cde0bea4c1b3e7c25f9a48f9bfb3296cf443c.\par
Rendering alone is not recorded as inspection. The main controller’s Chat/GPT-6/Pro UI observation is attributed to that controller; the host runtime self-description and lack of independent backend verification are recorded separately.\par
9 A. Giacchetto, R. Kramer, D. Lewański, and A. Sauvaget, The spin Gromov–Witten/Hurwitz correspondence for \(\mathbb P^1\), arXiv:2208.03259v3, 14 August 2025. https://arxiv.org/abs/2208.03259v3. EMS Press, journal metadata for DOI 10.4171/JEMS/1588. https://ems.press/journals/jems/articles/14298466. A. Giacchetto, R. Kramer, D. Lewański, and A. Sauvaget, Spin GW invariants of \(\mathbb P^1\): Calculations, nine-page author-linked PDF, retained as `spin\_GW\_calculations.pdf`; release date not independently established. Author’s papers page: https://reinierkramer.github.io/papers/.\par
\end{document}
